\documentclass[11pt,a4paper]{article}
\usepackage[T1]{fontenc}
\usepackage{lmodern}
\usepackage{microtype}
\usepackage[margin=2.3cm]{geometry}
\usepackage{amsmath,amssymb}
\usepackage{booktabs}
\usepackage{array}
\usepackage{enumitem}
\usepackage{xcolor}
\usepackage{listings}
\usepackage[colorlinks=true,linkcolor=blue!50!black,urlcolor=blue!50!black]{hyperref}

\lstset{language=Python, basicstyle=\ttfamily\footnotesize, keywordstyle=\color{blue!60!black},
        commentstyle=\color{green!40!black}, stringstyle=\color{red!50!black}, columns=fullflexible,
        keepspaces=true, frame=single, rulecolor=\color{black!25}, xleftmargin=4pt, xrightmargin=4pt,
        breaklines=true, showstringspaces=false}

\newcommand{\qp}{(q/p)}
\newcommand{\vb}[1]{\vec{#1}}
\newcommand{\br}[1]{\texttt{#1}}
\newcommand{\scal}{\sigma_{\mathrm{cal}}}
\newcommand{\mKKpi}{m_{KK\pi}}

\title{Per-candidate mass uncertainty $\sigma(\mKKpi)$\\ from momentum-calibrated tracks}
\author{R(D$_s$) analysis --- working note}
\date{3 November 2026}

\emergencystretch=2em
\begin{document}
\maketitle

\section{The idea}

The per-candidate uncertainty $\sigma(\mKKpi)$ is computed from what each track's own fit says about its
momentum, propagated linearly to the invariant mass, with the resolution added by the momentum calibration
folded in. The vertex fit plays no part. The result is a function of reconstructed quantities only, so it is
defined identically in data and simulation and can be used to build categories.

\begin{enumerate}[label=\textbf{\arabic*.}]
  \item \textbf{Inputs.} For each of $K_1$, $K_2$, $\pi$ take the $3\times3$ block of the helix covariance in
        $(\qp, \lambda, \phi)$. The three tracks are uncorrelated.
  \item \textbf{Propagation.} $\sigma^2(m) = \sum_i J_i\, C_i\, J_i^{T}$, with
        $J_i = \partial m/\partial(\qp,\lambda,\phi)_i$ evaluated at the calibrated momenta.
  \item \textbf{Calibration term.} Add to each track's $\qp$ variance the extra relative resolution that the
        calibration assigns to its $(p, p_T, \eta, q)$ bin, computed without the random number $z$.
  \item \textbf{Same definition everywhere.} Apply steps 1--3 unchanged to data and MC. Never use anything
        that exists only in MC ($z$, the realised smear, generator information).
\end{enumerate}

The muon is not part of $\mKKpi$ and never enters. Everything below generalises to $m_{KK}$ or $m_{KK\pi\mu}$
by changing the set of tracks in the sum.

\section{Step 1: the inputs}

Only the $3\times3$ momentum block of each track's covariance enters the mass, and the three tracks are
independent.

\paragraph{What the $5\times5$ matrix describes.} A CMS track is a helix with five parameters at its
reference point (the point of closest approach used by the packed candidate). They are stored, in this
order, as \br{<trk>\_cov\_<a>\_<b>} (upper triangle):

\begin{center}
\begin{tabular}{lll}
\toprule
parameter & meaning & unit \\
\midrule
$q/p$     & charge over total momentum & GeV$^{-1}$ \\
$\lambda$ & dip angle, $\lambda = \pi/2-\theta$; \ $\sin\lambda=\tanh\eta$, $\cos\lambda = 1/\cosh\eta$ & rad \\
$\phi$    & azimuth of the momentum at the reference point & rad \\
$d_{xy}$  & transverse impact parameter & cm \\
$d_{sz}$  & longitudinal impact parameter & cm \\
\bottomrule
\end{tabular}
\end{center}

The momentum depends on the first three only:
\begin{equation}
  \vb p = p\,(\cos\lambda\cos\phi,\; \cos\lambda\sin\phi,\; \sin\lambda), \qquad p = \frac{q}{\qp}.
  \label{eq:p}
\end{equation}

\paragraph{Why only the $3\times3$ block.} The invariant mass depends on the momenta, not on where the
tracks are: $\partial m/\partial d_{xy} = \partial m/\partial d_{sz} = 0$, so the last two columns of the
Jacobian (Step~2) vanish. The $\phi$--$d_{xy}$ and $\lambda$--$d_{sz}$ correlations are large in the
$5\times5$ matrix, but they multiply a zero and drop out. The correlations inside the block
($\qp$--$\lambda$, $\qp$--$\phi$, $\lambda$--$\phi$) do enter and must be kept.

This is a first-order statement. Tracks bend between the reference point and the D$_s$ vertex, so the
opening angles at the vertex differ by a few mrad from those at the reference point. That is a shift of the
mass value, identical in data and MC, not an uncertainty. It is the piece the vertex fit would have brought
in, together with its mismodelling.

\paragraph{Why the tracks are uncorrelated.} Each track is fitted from its own hits, whose measurement
errors are independent. The three $5\times5$ matrices are therefore independent blocks of one
$15\times15$ matrix, with zeros between them. Two effects break this, and neither belongs in a per-event
$\sigma$:
\begin{itemize}
  \item shared hits between close tracks: rare, and the $K^+K^-$ from a $\phi$ open up enough for this not
        to matter;
  \item common systematics (magnetic-field scale, alignment, material), which move all tracks of an event
        together: scale effects, corrected by the calibration, not random per-event spread.
\end{itemize}
After a vertex fit the refitted momenta become correlated, because they share one vertex; that is why the
fitted covariance has cross terms. Since we do not use the vertex fit, the covariance is block-diagonal and
the propagation is a sum over tracks.

\section{Step 2: propagation}

To first order each track contributes $J_i C_i J_i^T$, where $J_i$ is the gradient of the mass with
respect to that track's three helix parameters; the contributions add because the tracks are independent:
\begin{equation}
  \sigma^2(m) = \sum_{i\in\{K_1,K_2,\pi\}} J_i\, C_i\, J_i^{T}, \qquad
  J_i = \left( \frac{\partial m}{\partial \qp_i},\; \frac{\partial m}{\partial \lambda_i},\;
               \frac{\partial m}{\partial \phi_i} \right),
  \label{eq:prop}
\end{equation}
with $C_i$ the $3\times3$ block of Step~1. The chain rule splits $J_i$ into two factors: the mass with
respect to the momentum vector, and the momentum vector with respect to the helix parameters.

\paragraph{Mass with respect to one track's momentum.} With $E=\sum_i E_i$, $\vb P = \sum_i \vb p_i$ and
$E_i = \sqrt{p_i^2+m_i^2}$, differentiating $m^2 = E^2 - |\vb P|^2$ gives
\begin{equation}
  \vb g_i \equiv \frac{\partial m}{\partial \vb p_i}
  = \frac{1}{m}\left( E\,\frac{\vb p_i}{E_i} - \vb P \right)
  = \gamma\,\big( \vb\beta_i - \vb\beta \big),
  \label{eq:g}
\end{equation}
where $\vb\beta_i = \vb p_i/E_i$ is the velocity of track $i$, and $\vb\beta=\vb P/E$, $\gamma = E/m$ belong
to the $KK\pi$ system. The sensitivity of the mass to a track is the D$_s$ boost times the track's velocity
relative to the D$_s$. Two limits:
\begin{itemize}
  \item D$_s$ at rest: $\vb g_i = \vb\beta_i$. Only the momentum magnitude matters,
        $\partial m/\partial |\vb p_i| = \beta_i$, and a small rotation of one track leaves $m$ unchanged at
        first order.
  \item Boosted D$_s$: $\vb g_i$ acquires a large component transverse to the track, and direction errors
        start to matter (the opening-angle term).
\end{itemize}

\paragraph{Momentum vector with respect to the helix parameters.} For $|q|=1$, from $p = q/\qp$ and
Eq.~\eqref{eq:p},
\begin{equation}
  \frac{\partial \vb p}{\partial \qp} = -q\,p\,\vb p, \qquad
  \frac{\partial \vb p}{\partial \lambda} = p\,\hat e_\lambda, \qquad
  \frac{\partial \vb p}{\partial \phi} = p_T\,\hat e_\phi,
\end{equation}
with the unit vectors
\begin{equation}
  \hat e_\lambda = (-\sin\lambda\cos\phi,\; -\sin\lambda\sin\phi,\; \cos\lambda), \qquad
  \hat e_\phi = (-\sin\phi,\; \cos\phi,\; 0).
\end{equation}
(The first follows from $\mathrm{d}p/\mathrm{d}\qp = -q/\qp^2 = -q\,p^2$ and $\vb p = p\,\hat p$.)

\paragraph{The Jacobian.} Projecting each column on $\vb g_i$:
\begin{equation}
  J_i = \Big( -q_i\,p_i\,(\vb g_i\cdot\vb p_i),\;\; p_i\,(\vb g_i\cdot\hat e_{\lambda,i}),\;\;
              p_{T,i}\,(\vb g_i\cdot\hat e_{\phi,i}) \Big).
  \label{eq:J}
\end{equation}

\paragraph{Reading each term.} Since $\sigma\qp/|\qp| = \sigma_p/p$, the $\qp$ term of track $i$ alone gives
\begin{equation}
  \delta m_i\big|_{p} = (\vb g_i\cdot\hat p_i)\;\sigma_{p,i},
\end{equation}
the component of the relative velocity along the track times the momentum error. The two angle terms give
$p_i\,(\vb g_i\cdot\hat e_{\lambda,i})\,\sigma_{\lambda,i}$ and
$p_{T,i}\,(\vb g_i\cdot\hat e_{\phi,i})\,\sigma_{\phi,i}$. All three are combined through the full $C_i$,
correlations included, not simply added in quadrature.

\paragraph{Where to evaluate $J$.} Use the calibrated momenta (\br{pt}, \br{eta}, \br{phi}) with the kaon and
pion mass hypotheses, and the measured $\mKKpi$ in $1/m$: these are the momenta that define the mass being
categorised. The calibration does not change directions and rescales $|\vb p|$ by a factor $f$ close to 1,
so $J$ moves only at second order. Strictly, $C_i$ refers to the uncalibrated $\qp$, so its $\qp$ row and
column scale by $1/f$; a per-mille effect, but free to apply.

\section{Step 3: the calibration term}

The reco covariance describes the tracker's resolution as simulated. The calibration exists because the
resolution in data is worse, so the missing width is added to each track's $\qp$ variance. It is taken from
the calibration itself, as an expected width rather than the smear actually applied.

\paragraph{What the calibration does.} For every MC track, \br{TrackMomentumCalib} draws a standard
normal number and rescales the momentum, leaving the direction unchanged:
\begin{equation*}
  z = \mathrm{gauss}(\mathrm{run}, \mathrm{lumi}, \mathrm{event}, \mathrm{track\ key}), \qquad
  f = \frac{p'}{p} = \mathrm{factor}(p, p_T, \eta, q, z).
\end{equation*}
$f$ has two parts:
\begin{itemize}
  \item a scale part, its value at $z=0$, which shifts the momentum;
  \item a resolution part, its spread over $z$, which widens it.
\end{itemize}
Only the second is a resolution. The hadron terms apply (\br{kind='hadron'}); the muon terms are not needed
for $\mKKpi$.

\paragraph{The extra width, without $z$.} For the track's own $(p, p_T, \eta, q)$ define
\begin{equation}
  \sigma_{\mathrm{cal},i} = \frac{f(z=+1) - f(z=-1)}{2\, f(z=0)},
  \label{eq:scal}
\end{equation}
a relative width on $p$, and equally on $\qp$ since $\sigma\qp/|\qp| = \sigma_p/p$ at first order. If $f$
is not linear in $z$, use the RMS of $f/f(0)$ over $z\sim\mathcal N(0,1)$ instead; $f(\pm2)$ tells whether
that matters. Then
\begin{equation}
  C_i[\qp,\qp] \;\longrightarrow\; C_i[\qp,\qp] + \big(\sigma_{\mathrm{cal},i}\,\qp_i\big)^2 .
\end{equation}
In the mass this adds $\sigma_{\mathrm{cal},i}\; p_i\,(\vb g_i\cdot\hat p_i)$ in quadrature for each track: the momentum
term of Step~2 with $\sigma_p/p$ replaced by $\scal$.

\paragraph{Why in quadrature.} The calibration smear is drawn independently of the tracker measurement: it
uses its own $z$ and does not touch the direction, so it is uncorrelated with the angles. The $z$ of
different tracks are independent, so the term also adds track by track.

\paragraph{Why without $z$.} $\sigma_m$ must be a width: the expected mass error given the measured
kinematics, not the error that actually occurred. The realised $f(z)$ of a given MC track is MC truth, since
it says which way that track was smeared. Putting it into $\sigma_m$ would correlate $\sigma_m$ with the
actual mass residual in MC, which data could never reproduce. Data has no $z$ at all.

\paragraph{One assumption.} This presumes that the reco covariance already describes the resolution of
uncalibrated MC, so that adding $\scal$ gives the resolution of calibrated MC, hence of data. The MC pull
test of Section~\ref{sec:val} (check~2) tests exactly this. If it fails, the reco covariance needs its own
correction before anything is added.

\paragraph{Systematic variations.} Build the categories with the nominal calibration. The calibration
variations then move the resolution inside each category, a systematic on the per-category mass shapes; they
should not move the category boundaries.

\section{Step 4: one definition for data and MC}

$\sigma_m$ is a valid category variable only if it is the same function of reconstructed quantities in data
and MC. Then calibrated MC, which reproduces the kinematics and resolution of data, also reproduces how events
spread over $\sigma_m$ and the mass resolution inside each category.

\begin{center}
\begin{tabular}{>{\raggedright\arraybackslash}p{0.24\textwidth}>{\raggedright\arraybackslash}p{0.3\textwidth}>{\raggedright\arraybackslash}p{0.36\textwidth}}
\toprule
ingredient & data & MC \\
\midrule
helix covariance $C_i$ & packed candidate & packed candidate, same algorithm and format \\
momenta for $J_i$ & measured & calibrated (\br{pt}, \br{eta}, \br{phi}), not \br{pt\_raw} \\
$\sigma_{\mathrm{cal},i}$ & from the track's $(p,p_T,\eta,q)$, nominal & same function, same inputs \\
$z$, realised $f$, generator info & do not exist & exist, must not be used \\
\bottomrule
\end{tabular}
\end{center}

\paragraph{Why data gets $\scal$ too, although data is never smeared.} The tracker computes $C_i$ from its own
error model, identically for data and simulation. In MC that model describes the simulated resolution. In data
it describes the same simulated-quality resolution, which is smaller than the true one by exactly what the
calibration adds to MC. Hence $C_i\oplus\scal$ is the best estimate of the true width in both samples. Without
$\scal$ in data, its $\sigma_m$ would be biased low and the per-category widths would no longer match
calibrated MC.

\paragraph{What consistency buys, and what it does not.}
\begin{itemize}
  \item \textbf{Same meaning:} an event at a given $\sigma_m$ has the same expected resolution in data and MC.
        This follows from the definition, provided the assumption of Step~3 holds.
  \item \textbf{Same population:} the $\sigma_m$ distribution, hence the category fractions, agree between data
        and MC. This does \emph{not} follow from the definition: it depends on how well MC models the track
        kinematics and the $C_i$ themselves (hit patterns, $\eta$ and $p_T$ spectra), and must be measured
        (Section~\ref{sec:val}).
\end{itemize}

\paragraph{Ways to break it without noticing.}
\begin{itemize}
  \item using the vertex-fit covariance;
  \item using $z$ or \br{calib\_f} in the $\scal$ term;
  \item using \br{pt\_raw} in MC but measured momenta in data (different $p$ in $J$ and in the calibration bin);
  \item computing $\scal$ in MC only;
  \item a covariance rescaling applied to one sample only: the branches \br{<trk>\_cov\_scale\_*} record any
        \br{--cov-scale} factors; check that the stored covariance is the unscaled one, or divide the factors
        out, before using it in either sample.
\end{itemize}

\paragraph{A cheaper variant.} Leaving $\scal$ out in both samples is still a consistent definition, only a
worse estimator: its pulls are wider than 1. If the categories are quantile cuts and $\scal$ is a roughly
constant fraction of the total, the ranking hardly changes. Where $\scal$ varies strongly between bins (barrel
vs endcap, low vs high $p$) it reorders events, so including it is the better default.

\section{Validation}
\label{sec:val}

Five checks, in order, each testing a different assumption above. The first two can run on the current,
uncalibrated, signal MC.

\begin{enumerate}[label=\textbf{\arabic*.}]
  \item \textbf{Jacobian, numerically.} For a few hundred candidates move each helix parameter by a small step,
        recompute $\mKKpi$, and compare the finite difference with Eq.~\eqref{eq:J}. This catches sign and
        convention errors in the $\qp$ column and in $\lambda$ vs $\theta$.
  \item \textbf{MC pull, uncalibrated sample, without $\scal$.} On truth-matched D$_s\to\phi\pi$, form
        $(\mKKpi - m_{\mathrm{gen}})/\sigma_m$.
        \begin{itemize}
          \item Compute $m_{\mathrm{gen}}$ from the generator momenta of the matched tracks
                (branches \br{gen\_vis\_\{kp,km,pi\}\_*}): this removes FSR from the comparison; the D$_s$ natural width is
                negligible.
          \item Measure the core width (Gaussian fit to the core, or IQR$/1.349$): radiative tails are not part
                of $\sigma_m$.
          \item Expect width 1, flat in quantiles of $\sigma_m$, in $\eta$ and in $p$. This tests the assumption
                of Step~3.
          \item A width different from 1 but constant is a global scale factor; one that varies across
                $\sigma_m$ quantiles points to an error in $C_i$ or $J_i$.
        \end{itemize}
  \item \textbf{MC pull, calibrated sample, with $\scal$.} Same test, width 1 expected. This checks the
        quadrature sum of Step~3 against the smear actually applied.
  \item \textbf{D$_s$ peak in data vs calibrated MC, per category} --- the test that matters for the analysis.
        \begin{itemize}
          \item Split the candidates in quantiles of $\sigma_m$ (the categories, or finer).
          \item In each, fit the $\mKKpi$ peak with a free signal width and a linear or exponential background.
          \item Plot the fitted width against $\langle\sigma_m\rangle$ in the bin: data and calibrated MC should
                lie on the same line, ideally of slope 1 through the origin.
          \item Quote the per-category data/MC width ratio. It should be 1 within its error; otherwise it
                becomes a per-category resolution correction.
        \end{itemize}
  \item \textbf{Population.} Compare the $\sigma_m$ distribution of D$_s$ signal in data (sideband-subtracted or
        sPlot) with calibrated MC. The category fractions must agree. A residual difference means the MC
        kinematics or covariance modelling is off: reweight, or assign a systematic on the category yields.
        It says nothing about the $\sigma_m$ definition.
\end{enumerate}

Checks 4 and 5 use the D$_s\to\phi\pi$ peak itself, abundant in data and the same particle as in the signal:
no control channel is needed.

\section{Implementation with the existing branches}

Everything needed is in the current ntuples except the calibration applier. A numpy sketch, vectorised over
candidates:

\begin{lstlisting}
import numpy as np

MASS = {'k1': 0.493677, 'k2': 0.493677, 'pi': 0.13957039}
PARS = ('qoverp', 'lambda', 'phi')          # first three helix parameters

def cov3(d, t):
    '''3x3 (q/p, lambda, phi) block; the branches are the upper triangle.'''
    n = len(d[t + '_pt'])
    C = np.empty((n, 3, 3))
    for a in range(3):
        for b in range(a, 3):
            C[:, a, b] = C[:, b, a] = d['%s_cov_%s_%s' % (t, PARS[a], PARS[b])]
    f = np.nan_to_num(d[t + '_calib_f'].astype(float), nan=1.0)   # data / uncalibrated MC: 1
    C[:, 0, :] /= f[:, None]                # q/p' = (q/p)/f
    C[:, :, 0] /= f[:, None]
    return C

def sigma_m_kkpi(d, sigma_cal=None):
    '''m(KKpi) and its per-candidate uncertainty. sigma_cal: {track: relative
    calibration width per candidate}, or None to leave the term out.'''
    trk = {}
    for t, mt in MASS.items():
        pt, eta, phi = (d['%s_%s' % (t, q)].astype(float) for q in ('pt', 'eta', 'phi'))
        p3 = np.stack([pt * np.cos(phi), pt * np.sin(phi), pt * np.sinh(eta)], 1)
        p = pt * np.cosh(eta)
        trk[t] = (p3, np.sqrt(p ** 2 + mt ** 2), p, pt, eta, phi)
    P = sum(v[0] for v in trk.values())
    E = sum(v[1] for v in trk.values())
    m = np.sqrt(E ** 2 - np.sum(P ** 2, 1))
    var = np.zeros(len(m))
    for t, (p3, Ei, p, pt, eta, phi) in trk.items():
        g = (E[:, None] * p3 / Ei[:, None] - P) / m[:, None]          # dm/dp_vec
        q = d[t + '_charge'].astype(float)
        sl, cl = np.tanh(eta), 1.0 / np.cosh(eta)                       # sin, cos lambda
        e_lam = np.stack([-sl * np.cos(phi), -sl * np.sin(phi), cl], 1)
        e_phi = np.stack([-np.sin(phi), np.cos(phi), np.zeros_like(phi)], 1)
        J = np.stack([-q * p * np.sum(g * p3, 1),                      # d/d(q/p)
                      p * np.sum(g * e_lam, 1),                         # d/d(lambda)
                      pt * np.sum(g * e_phi, 1)], 1)                    # d/d(phi)
        C = cov3(d, t)
        if sigma_cal is not None:
            C[:, 0, 0] += (sigma_cal[t] / p) ** 2                       # (sigma_cal * q/p)^2
        var += np.einsum('ni,nij,nj->n', J, C, J)
    return m, np.sqrt(var)
\end{lstlisting}

The calibration width, once \br{rds\_trkcalib.py} and its JSON are available (\br{TrackCalib.factor} as
called in \br{TrackMomentumCalib.z\_and\_f}; wrap it in \br{np.vectorize} if it takes scalars):

\begin{lstlisting}
def sigma_cal_rel(cal, p, pt, eta, q, kind='hadron', variation='nominal'):
    f = lambda z: cal.factor(p, pt, eta, q, z, variation, kind=kind)
    return (f(1.0) - f(-1.0)) / (2.0 * f(0.0))
\end{lstlisting}

\paragraph{Notes.}
\begin{itemize}
  \item Feed the calibrated momenta to the bin lookup in MC and the measured ones in data. The calibration was
        defined on uncalibrated momenta, but $f$ is within a per mille of 1, so this matters only at bin edges,
        and it keeps the definition symmetric.
  \item Skip or flag candidates with \br{<trk>\_cov\_pos\_def} false; check what \br{<trk>\_cov\_fixed} means
        for them before trusting their covariance.
  \item Packed-candidate covariances are stored with reduced precision: fine for a per-event $\sigma$, and
        check~2 would show it otherwise.
  \item For the next production the same function can fill a convenience branch (\br{kkpi\_mass\_err}), so the
        analysis does not need the applier.
\end{itemize}

\section{Check yourself}

\begin{enumerate}[label=\textbf{\arabic*.}]
  \item A D$_s$ is produced at rest and decays to $K^+K^-\pi^+$. Which of the three helix-parameter errors
        drive $\sigma_m$, and why do the other two drop out at first order?
  \item Boost the same D$_s$ to 15\,GeV. Using $\vb g_i = \gamma(\vb\beta_i-\vb\beta)$, predict which tracks
        dominate $\sigma_m$: the slow or the fast ones, those at large or small angle to the D$_s$?
  \item The $\phi$--$d_{xy}$ correlation in the $5\times5$ covariance is strong for low-$p_T$ tracks. Does it
        change $\sigma_m$? What would change if $\mKKpi$ were computed with the momenta at the D$_s$ vertex
        instead?
  \item A colleague adds $\scal$ to MC only, ``because only MC was smeared''. How do the per-category D$_s$
        peak widths in data and MC now disagree, and in which direction?
  \item In check~2 the pull width is 1.15 in the lowest $\sigma_m$ quantile and 1.00 in the highest. Which
        ingredient do you suspect first, and which plot tells you?
\end{enumerate}

\end{document}
